\section{Two oscillators and finite energy}\label{sec:oscillators}
All logarithms in entanglement entropies are base two; trace entropies and modular variables below use natural logarithms. Oscillator energies are occupation numbers, with the zero point omitted. Arbitrary bounded local measurements are allowed. In particular, these statements impose no Gaussian restriction or finite control-time bound.

\begin{theorem}[An explicit finite-energy perfect strategy]\label{thm:oscillators}
The perfect strategy of~\eqref{eq:spatialstate} is a strategy on two harmonic oscillators with
\[
 \langle N_A\rangle=\langle N_B\rangle=\frac{275}{18},\qquad
 \langle N_A+N_B\rangle=\frac{275}{9}.
\]
Every polynomial number moment is finite. These are the least possible occupations among perfect pure spatial strategies on two oscillators, and the passive minima for the displayed Schmidt spectrum. Retaining $K$ levels gives the loss~\eqref{eq:exactloss} and mean occupation per side
\[
 \langle N\rangle_K=\frac{275}{18}-\frac{20K}{4^K-1}.
\]
\end{theorem}
\begin{proof}
Identify $A(i,j,k)$ with $|20k+4j+i\rangle$ and $B(j,k)$ with $|20k+16+j\rangle$. This is a bijection onto the oscillator basis. The probabilities, in decreasing order, are sixteen copies of $4^{-k}/24$ and then four of $4^{-k}/48$, at each $k\ge0$. They are decreasing even across adjacent levels. Their level mass is $(3/4)4^{-k}$, with expected level $1/3$. The unnormalized mean offset in a level is
\[
 \frac{0+\cdots+15}{24}+\frac{16+\cdots+19}{48}=\frac{155}{24}.
\]
Hence the mean is $20/3+(155/24)(4/3)=275/18$. Decreasing probabilities assigned to increasing energies minimize their mean: finite-dimensional rearrangement, followed by spectral cutoffs and monotone convergence, gives this passive lower bound for every placement. The displayed placement attains it. Global optimality over perfect pure strategies follows from the integer-spectrum argument in Theorem~\ref{thm:energy-minimum}.

The generating function is
\[
 \langle z^N\rangle=
 \frac{\sum_{r=0}^{15}z^r/24+\sum_{r=16}^{19}z^r/48}{1-z^{20}/4}.
\]
Its radius exceeds one, giving every polynomial moment. The local measurements are the orthogonal direct sums of the finite ideal PVMs already described in Section~\ref{sec:strategy}, plus the boundary and label projections. They are everywhere-defined projections of norm at most one, and their finite outcome sums are the identity in the strong topology. Thus the oscillator identification introduces no unbounded measurement-domain assumption. The block and link proof of Theorem~\ref{thm:main} verifies all fifty-two pairs. Finally, conditioning the geometric level distribution on $k<K$ gives mean level $1/3-K/(4^K-1)$; the offset distribution is unchanged, proving the truncation formula.
\end{proof}

\subsection{Attainment of bounded resource prices}
For a game $G$, let $S_G(\eps)$ be the infimum of pure-state entanglement entropy among spatial strategies of loss at most $\eps$. Let $E_G(\eps)$ be the corresponding total mean energy on a specified pair of nonnegative local Hamiltonians, allowing normal mixed states and arbitrary POVMs. Empty infima are infinite. These definitions specify the state domain; a finite-dimensional infimum at loss zero would be a different quantity.

\begin{theorem}[Three compactness criteria]\label{thm:attainment}
For a finite game, bounded dimension as loss tends to zero is equivalent to finite-dimensional perfect attainment. A finite limit $s=\lim_{\eps\downarrow0}S_G(\eps)$ is attained by a perfect spatial pure strategy of entropy $s$. If both local Hamiltonians have compact resolvent, a finite limit $e=\lim_{\eps\downarrow0}E_G(\eps)$ is attained by a perfect strategy, and a perfect pure strategy has energy $e$ as well.
\end{theorem}
\begin{proof}
In a fixed finite dimension, density operators and the finitely many POVMs form compact sets and the loss is continuous. A convergent subsequence proves the dimension statement.

For entropy choose pure strategies with losses tending to zero and entropies tending to $s$. Transport their Schmidt supports to one common countable basis, with decreasing probabilities $p_i^{(n)}$. Since $p_i^{(n)}\le1/i$, their tails past $r>1$ are at most $(s+1)/\log_2r$, after discarding a finite prefix. Coordinatewise subselection and this tightness give convergence in $\ell^1$ to a probability vector $p$. The canonical pure vectors converge in norm because $(\sqrt x-\sqrt y)^2\le|x-y|$. Extract simultaneous ultraweakly convergent subsequences of all local effects. Positivity and completeness persist. Tensor products converge weakly: verify matrix elements on product vectors, then finite sums, and use uniform operator bounds and density. Combining this with norm convergence of vectors makes the limiting loss zero. Fatou gives $S(p)\le s$. Every perfect strategy bounds every approximate price from above, so equality holds. Transport to Schmidt supports avoids an unjustified extension of an infinite-dimensional basis map to a unitary on its original kernel.

For energy put $H=h_A\otimes I+I\otimes h_B$. Its bounded-energy spectral projections have finite rank. If $\operatorname{tr}(\rho H)\le e+1$, then $\operatorname{tr}(\rho1_{(R,\infty)}(H))\le(e+1)/R$. Finite-rank compression and the square-root tail estimate give trace-norm precompactness of these density operators. Extract a trace-norm limit and weak limits of the measurements as above. The limit is perfect, and lower semicontinuity of energy, obtained by bounded spectral cutoffs, gives energy at most $e$; it must equal $e$. In a pure decomposition of this state, positivity of the loss makes every positive-weight component perfect. Each has energy at least $e$, so some, indeed almost every component of finite weight, has energy $e$. This proves the pure assertion without assuming the minimizing state was pure.
\end{proof}

These criteria separate attainment from computation. For $\Gamma$, the exact entropy and one-sided oscillator endpoints are computed in Theorem~\ref{thm:price-limits}, using the spectral classification of Section~\ref{sec:spectra}.

\begin{proposition}[A comparison with matched logarithmic games]\label{prop:oscillator-exponent}
Suppose a rational value-one game has matched finite-dimensional laws with $\alpha>0$:
\[
 \log_2D_G(\eps)=\alpha\log_2(1/\eps)+o(\log(1/\eps)),\qquad
 d_G^\to(\eps)=\alpha\log_2(1/\eps)+o(\log(1/\eps)),
\]
where $d_G^\to$ is the original-state one-way distillable-entanglement demand, and maximally entangled states attain the dimension upper law. With $m$ oscillator modes per party its optimal energy has exponent $\alpha/m$:
\[
 E_G(\eps)=\eps^{-\alpha/m+o(1)}.
\]
\end{proposition}
\begin{proof}
The maximum entropy of $m$ oscillators at total mean occupation $E$ is $m g(E/m)$, where $g(x)=(x+1)\log_2(x+1)-x\log_2x$. Relative entropy against the product thermal state proves this bound, and $g(x)\le\log_2(\mathrm e(x+1))$. Since directional distillation is at most either marginal entropy, each mean occupation is at least $m(2^{d_G^\to(\eps)/m}/\mathrm e-1)$. For normal spatial strategies, finite local truncations approximate their correlations with arbitrarily small relative loss slack; local monotonicity of distillation gives the same lower bound with $\eps$ replaced by $(1+o(1))\eps$.

A rank-$D$ upper strategy fits into the $m$-mode grid with each coordinate less than $\lceil D^{1/m}\rceil$. Its mean occupation is at most $m\lceil D^{1/m}\rceil$. The two laws squeeze the exponent. The matched games in~\cite{DMprice} supply examples. This proposition uses both matched laws; entropy divergence alone does not imply this power law, and $\Gamma$ has the separate bounded-energy behavior above.
\end{proof}

\subsection{Global oscillator optima and small-loss limits}\label{sec:energy-optima}
Let $\Lambda_*=(p_0^*,p_1^*,\ldots)$ be the decreasing probabilities of Theorem~\ref{thm:oscillators}: sixteen copies of $4^{-k}/24$, then four of $4^{-k}/48$, at every level. Section~\ref{sec:spectra} proves that this is the unique majorization maximum among perfect type-I marginal spectra.

\begin{theorem}[The least oscillator occupation]\label{thm:energy-minimum}
For any perfect strategy with a normal type-I Alice marginal and any nonnegative compact-resolvent local Hamiltonian with ordered eigenvalues $e_0\le e_1\le\cdots$,
\[
 \langle H_A\rangle\ge\sum_{i\ge0}e_i p_i^*.
\]
When the series is finite it is attained by preparing the displayed strategy in the ordered energy eigenbasis. If the series is infinite, every perfect strategy has infinite energy on that platform. In particular $\langle N_A\rangle\ge275/18$. Equality for the oscillator requires spectrum $\Lambda_*$ and passive placement. For pure spatial strategies, each side has mean occupation at least $275/18$, and the total at least $275/9$, all attained by the displayed strategy.
\end{theorem}
\begin{proof}
For a rank-$r$ projection, the variational principle bounds its probability by the sum of the first $r$ marginal eigenvalues. Expand the Hamiltonian as a monotone sum of positive steps at its energy-prefix projections. With $p$ the ordered marginal spectrum this gives
\[
 \langle H_A\rangle\ge e_0+\sum_{r\ge1}(e_r-e_{r-1})T_r(p)
          \ge e_0+\sum_{r\ge1}(e_r-e_{r-1})T_r(\Lambda_*)
          =\sum_i e_i p_i^*.
\]
Theorem~\ref{thm:majorization} proves the second inequality. All sums are nonnegative, so no finite-expectation assumption is needed. Ordered preparation attains the expression. For $N_A$ every increment is one. Equality forces every spectral tail to equal the released tail, hence $p=\Lambda_*$; equality in the energy-prefix inequalities is passive placement. The computed value is $275/18$. Pure strategies have the same nonzero Schmidt spectrum on Bob, giving the other side and the total.
\end{proof}

Attainment here permits preparation and transport on the Schmidt support. It is not an assertion that every entire-space unitary orbit reaches a passive minimum: a prescribed marginal with an extra kernel retains it on its unitary orbit, while the oscillator's full-support passive optimum has none. Nor does compact resolvent imply a finite bound; for instance $e_i=2^i-1$ makes the displayed series infinite. The pure Bob-side conclusion is the only Bob-side energy assertion used here.

Define $E_{A,\Gamma}(\eps)$ as the infimum of Alice's mean oscillator occupation over normal spatial strategies of loss at most $\eps$, allowing mixed states and unrestricted Bob dimension and POVMs. The entropy price $S_\Gamma$ retains the pure spatial convention above. Finite-only prices at positive loss are defined by the same infima restricted to finite-dimensional strategies.

\begin{samepage}
\begin{theorem}[Exact limiting resource prices]\label{thm:price-limits}
\[
 \lim_{\eps\downarrow0}S_\Gamma(\eps)=\log_2 24+7/9,
 \qquad \lim_{\eps\downarrow0}E_{A,\Gamma}(\eps)=275/18.
\]
The same limits hold for the finite-only prices at positive loss. These are limits, not quantitative convergence rates; the finite-only feasible set at zero loss is empty.
\end{theorem}
\end{samepage}
\begin{proof}
Monotonicity and the perfect released strategy give finite upper limits. The entropy compactness argument in Theorem~\ref{thm:attainment} produces a perfect pure limit at the limiting entropy. Theorem~\ref{thm:renyi-minimum} bounds that entropy below by $H_\Gamma$, so equality follows.

For one-sided energy, purify any mixed strategy into Bob without changing Alice's marginal, cost or correlations. Order its Schmidt probabilities $p_i$ and transport both Schmidt supports to the canonical vector $\sum_{i\ge0}\sqrt{p_i}|ii\rangle$, placing Alice passively. Her new mean $\sum_i i p_i$ is no greater than the original energy, and Bob is unconstrained. Along a sequence approaching a finite limiting cost $e$,
\[
 \sum_{i\ge r}p_i\le(e+o(1))/r.
\]
Coordinatewise subselection and this tightness give $\ell^1$ convergence of the lists and norm convergence of canonical pure vectors. Simultaneous ultraweak subsequences of the finitely many local effects remain POVMs. Product effects converge weakly first on product vectors and then by density, so the limiting loss is zero. Fatou bounds its Alice occupation by $e$. Theorem~\ref{thm:energy-minimum} forces $e\ge275/18$. This controls Bob by Schmidt transport, without applying two-sided Hamiltonian compactness to an unconstrained Bob system.

For finite-only prices the original $K$-level compressions have loss $1/[624(4^K-1)]$, energy tending to $275/18$ and entropy
\[
 H_K=H_\Gamma-\frac{2K}{4^K-1}+\log_2(1-4^{-K})\longrightarrow H_\Gamma.
\]
They supply the upper limits. Restriction to a smaller feasible class supplies the lower limits already proved.
\end{proof}

